\chapter{Black--Scholes Three Ways}\label{ch:dv:black-scholes-three-ways}

On 26 April 1973 the Chicago Board Options Exchange opened for business with
calls on 16 stocks and traded 911 contracts. In the same spring the
\emph{Journal of Political Economy} printed a paper by Fischer Black and Myron
Scholes, with a companion by Robert Merton, that gave the price of such a call
as a closed formula in five numbers: spot, strike, time, rate and
volatility. The formula was right for a reason that surprised economists: the
expected return of the share, the one number every investor argues about, is
not among the five. This chapter derives the formula three times, because each
derivation carries a different lesson: replication shows why the drift
disappears and what a hedger actually does; the martingale route shows which
expectation is being taken and how to generalise it; the binomial limit shows
that nothing here needs continuous-time mathematics to be true, only to be
convenient.

\section{The model and the replication argument}

\begin{definition}[Black--Scholes model]\label{def:dv:black-scholes-three-ways:model}
The \emph{Black--Scholes model}\index{Black--Scholes model} is a market with a
bank account growing at a constant rate $r$ and a share that pays a continuous
dividend yield $q$ and follows a geometric Brownian motion (One Quant Book 4,
chapter 4) under the real-world measure,
\[
dS_t=\mu S_t\,dt+\sigma S_t\,dW_t,
\]
with constant $\mu$ and $\sigma>0$. Trading is continuous and frictionless,
short sales are allowed, and cash can be borrowed and lent at $r$.
\end{definition}

Let $V(t,S)$ be the value of a European claim paying $g(S_T)$, assumed smooth in
$(t,S)$. Hold the claim and sell $\Delta$ shares. The short position pays the
dividend $qS\,dt$ to the lender of the share. By Itô's formula (One Quant Book 4,
chapter 3) the position $\Pi=V-\Delta S$ changes by
\[
d\Pi=\Bigl(\partial_tV+\tfrac12\sigma^2S^2\partial_{SS}V-q\Delta S\Bigr)dt
+\bigl(\partial_SV-\Delta\bigr)\,dS .
\]
Choosing $\Delta=\partial_SV$ removes the only random term. A portfolio with no
risk must earn the riskless rate, or borrowing to buy it (or selling it to
lend) would be an arbitrage: $d\Pi=r\Pi\,dt=r(V-S\partial_SV)\,dt$. Equating the
two expressions of $d\Pi$ gives the pricing equation.

\begin{definition}[Black--Scholes equation]\label{def:dv:black-scholes-three-ways:pde}
The \emph{Black--Scholes equation}\index{Black--Scholes equation} is the
parabolic partial differential equation
\[
\partial_tV+(r-q)S\,\partial_SV+\tfrac12\sigma^2S^2\,\partial_{SS}V-rV=0,
\qquad V(T,S)=g(S),
\]
satisfied by the value of any European claim on the share in the \omterm{def:dv:black-scholes-three-ways:model}{Black--Scholes
model}.
\end{definition}

The drift $\mu$ multiplied $dS$, and $dS$ was hedged away: it cancelled with the
position in shares. What survived is the variance term $\tfrac12\sigma^2S^2
\partial_{SS}V$, which no position in shares can cancel because it comes from the
curvature of $V$, and the financing of the hedge at $r-q$. The equation is
linear, so the value of a portfolio of claims is the sum of the values, and it
is the same for every investor whatever their view of $\mu$.

\begin{remark}[What the hedger does]\label{rem:dv:black-scholes-three-ways:hedger}
Selling the claim and holding $\Delta=\partial_SV$ shares financed at $r$ is a
\omterm{def:dv:no-arbitrage-and-the-fundamental-theorems:sf}{self-financing strategy} (chapter 1) whose value tracks $V$ exactly in the model.
\Cref{fig:dv:black-scholes-three-ways:replication} runs it once, rebalanced
daily along a path simulated with a drift of 12\% a year: the portfolio follows
the option, and the gap at expiry, 0.39 on a premium of 10.45, is the cost of
hedging daily instead of continuously, the subject of chapter 4.
\end{remark}

\begin{omfigure}
\begin{tikzpicture}
\begin{axis}[width=11cm,height=5cm,
  xlabel={time (years)},ylabel={value},
  tick label style={font=\small},label style={font=\small},
  xmin=0,xmax=1,ymin=0,ymax=26,grid=major,grid style={gray!25},
  legend columns=2,legend style={font=\footnotesize,at={(0.5,-0.3)},anchor=north,draw=none,fill=none,/tikz/every even column/.append style={column sep=8pt}},
  table/col sep=comma]
\addplot[omDef,very thick] table[x=t,y=option]{figdata/derivatives/03-black-scholes-three-ways/replication.csv};
\addplot[omThm,thick,dashed] table[x=t,y=portfolio]{figdata/derivatives/03-black-scholes-three-ways/replication.csv};
\legend{call value by the formula,replicating portfolio (daily rebalancing)}
\end{axis}
\end{tikzpicture}

\omcaption{Replication along one path. A one-year at-the-money call ($r=5\%$,
$\sigma=20\%$) is replicated by holding $\Delta$ shares and borrowing, rebalanced
252 times; the share, simulated with a drift of 12\%, ends at 115.46. The two
curves are indistinguishable at this scale; the final gap is 0.39. Data: the tutorial.}\label{fig:dv:black-scholes-three-ways:replication}
\end{omfigure}

\section{The pricing equation}

The \omterm{def:dv:black-scholes-three-ways:pde}{Black--Scholes equation} is a backward heat equation in disguise. With
$x=\ln S$ and time to expiry $\tau=T-t$, and $V=e^{-r\tau}u$, it becomes
$\partial_\tau u=(r-q-\tfrac12\sigma^2)\partial_xu+\tfrac12\sigma^2\partial_{xx}u$, a
diffusion with constant coefficients. By the Feynman--Kac formula (One Quant
Book 4, chapter 4) its solution is an expectation:
\begin{equation}\label{eq:dv:black-scholes-three-ways:fk}
V(t,S)=e^{-r(T-t)}\,\E\bigl[g(S_T)\mid S_t=S\bigr],\qquad
dS_u=(r-q)S_u\,du+\sigma S_u\,d\widetilde W_u .
\end{equation}
The expectation is taken for a share that grows at $r-q$, not at $\mu$: the
equation knows nothing of $\mu$, so neither does its solution. This is where the
phrase ``risk-neutral'' comes from: prices are expectations as if investors
required no premium for risk. It is a statement about prices, which is how
chapter 1 read it, not about investors.

\begin{omfigure}
\begin{tikzpicture}
\begin{axis}[width=10.5cm,height=5cm,
  xlabel={spot},ylabel={call value},
  tick label style={font=\small},label style={font=\small},
  xmin=60,xmax=140,ymin=0,ymax=45,grid=major,grid style={gray!25},
  legend columns=4,legend style={font=\footnotesize,at={(0.5,-0.3)},anchor=north,draw=none,fill=none,/tikz/every even column/.append style={column sep=8pt}},
  table/col sep=comma]
\addplot[gray,thick,dashed] table[x=spot,y=payoff]{figdata/derivatives/03-black-scholes-three-ways/call_vs_spot.csv};
\addplot[omProp,thick] table[x=spot,y=t01]{figdata/derivatives/03-black-scholes-three-ways/call_vs_spot.csv};
\addplot[omDef,thick] table[x=spot,y=t05]{figdata/derivatives/03-black-scholes-three-ways/call_vs_spot.csv};
\addplot[omThm,very thick] table[x=spot,y=t1]{figdata/derivatives/03-black-scholes-three-ways/call_vs_spot.csv};
\legend{payoff,0.1 year,0.5 year,1 year}
\end{axis}
\end{tikzpicture}

\omcaption{The call struck at 100 as a solution of the pricing equation, at
three times to expiry ($r=5\%$, $\sigma=20\%$). Going back in time the payoff's
kink at the strike is smoothed out by diffusion and the curve is lifted by the
interest on the strike. Data: the chapter's code.}\label{fig:dv:black-scholes-three-ways:values}
\end{omfigure}

\section{The martingale route}

Chapter 1 priced a claim as a discounted expectation under a measure in which
discounted traded prices are martingales. In the \omterm{def:dv:black-scholes-three-ways:model}{Black--Scholes model} Girsanov's
theorem (One Quant Book 4, chapter 5) gives it explicitly: under the
risk-neutral measure $\mathbb Q$ the process $W^{\mathbb Q}_t=W_t+\frac{\mu-r+q}
{\sigma}t$ is a Brownian motion and
\[
S_T=S_0\exp\Bigl(\bigl(r-q-\tfrac12\sigma^2\bigr)T+\sigma W^{\mathbb Q}_T\Bigr),
\]
so that $e^{-(r-q)t}S_t$ is a $\mathbb Q$-martingale: the dividend-adjusted share,
discounted at $r$, is a fair game. The price of the call is
\[
C=e^{-rT}\E^{\mathbb Q}\bigl[(S_T-K)\mathbf 1_{S_T>K}\bigr]
=e^{-rT}\E^{\mathbb Q}\bigl[S_T\mathbf 1_{S_T>K}\bigr]-Ke^{-rT}\,\mathbb Q(S_T>K).
\]
The second term is a normal probability. The first needs one more idea: change
the numeraire to the share (with its dividends reinvested). Under the share
measure $\mathbb Q^S$, with density $S_Te^{qT}/(S_0e^{rT})$, the first expectation
becomes $S_0e^{(r-q)T}\,\mathbb Q^S(S_T>K)$, and under $\mathbb Q^S$ the log-share
has drift $r-q+\tfrac12\sigma^2$ instead of $r-q-\tfrac12\sigma^2$.

\begin{theorem}[Black--Scholes formula]\label{thm:dv:black-scholes-three-ways:formula}
In the \omterm{def:dv:black-scholes-three-ways:model}{Black--Scholes model} the European call and put are worth
\[
C=Se^{-qT}\Phi(d_1)-Ke^{-rT}\Phi(d_2),\qquad
P=Ke^{-rT}\Phi(-d_2)-Se^{-qT}\Phi(-d_1),
\]
\[
d_{1,2}=\frac{\ln(S/K)+(r-q)T}{\sigma\sqrt T}\pm\tfrac12\sigma\sqrt T,
\]
the \emph{Black--Scholes formula}\index{Black--Scholes formula}. Here
$\Phi(d_2)=\mathbb Q(S_T>K)$ is the risk-neutral probability of exercise and
$\Phi(d_1)=\mathbb Q^S(S_T>K)$ the same probability under the share measure, and
$e^{-qT}\Phi(d_1)=\partial_SC$ is the hedge ratio.
\end{theorem}
\begin{proof}
$\ln S_T$ is normal with variance $\sigma^2T$ under both measures and with means
differing by $\sigma^2T$; $\mathbb Q(S_T>K)=\Phi(d_2)$ and $\mathbb Q^S(S_T>K)=
\Phi(d_1)$ follow by standardising. The put follows by put--call parity. The
identity $\partial_SC=e^{-qT}\Phi(d_1)$ uses $S e^{-qT}\varphi(d_1)=Ke^{-rT}
\varphi(d_2)$, which cancels the terms from differentiating $d_1$ and $d_2$.
\end{proof}

\begin{example}[The desk's reference call]\label{ex:dv:black-scholes-three-ways:ref}
$S=K=100$, $T=1$, $r=5\%$, $q=0$, $\sigma=20\%$: $d_1=0.35$, $d_2=0.15$,
$\Phi(d_1)=0.6368$, $\Phi(d_2)=0.5596$, $C=10.4506$, $P=5.5735$. The risk-neutral
probability of exercise is 0.56; with a real-world drift of 12\% the
probability that the call ends in the money is 0.69, and it plays no role in the
price.
\end{example}

\begin{proposition}[At-the-money approximation]\label{prop:dv:black-scholes-three-ways:atm}
At the forward, $K=F=Se^{(r-q)T}$, the call and the put are both worth
$e^{-rT}F\bigl(2\Phi(\tfrac12\sigma\sqrt T)-1\bigr)\approx0.4\,e^{-rT}F\sigma\sqrt T$.
\end{proposition}
\begin{proof}
With $K=F$, $d_{1,2}=\pm\tfrac12\sigma\sqrt T$ and $C=e^{-rT}F(\Phi(d_1)-\Phi(-d_1))$;
expand $\Phi$ at 0 with slope $\varphi(0)=1/\sqrt{2\pi}\approx0.3989$.
\end{proof}

\section{The binomial limit}

The \omterm{def:dv:the-binomial-model:crr}{Cox--Ross--Rubinstein tree} of chapter 2 needs no stochastic calculus and
reaches the same number. With $n$ steps the tree call is the discounted
expectation of $(S_T-K)^+$ under a binomial law of the number of up moves. As
$n\to\infty$ with $u=e^{\sigma\sqrt{T/n}}$, the tree's probability of finishing
above $K$ converges to $\Phi(d_2)$, and the same probability computed with the
weights $pu/R$ and $(1-p)d/R$ (the share measure on the tree) converges to
$\Phi(d_1)$: the two terms of the formula appear as limits of the two binomial
sums of \cref{prop:dv:the-binomial-model:formula}. For the reference call the
1\,000-step tree gives 10.4486, 0.0020 below the formula, the error of order
$1/n$ measured in chapter 2.

\section{Black's formula, and what each derivation teaches}

Many underlyings are quoted as forwards or futures: commodity futures, bond
futures, forward interest rates, FX forwards. Pricing on the forward removes the
dividend yield, the borrow and the financing from the formula; they are all
inside $F$.

\begin{definition}[Black model]\label{def:dv:black-scholes-three-ways:black}
The \emph{Black model}\index{Black model} assumes that the forward (or futures)
price $F_t$ of the underlying for delivery at $T$ is a driftless geometric
Brownian motion under the $T$-forward measure (One Quant Book 4, chapter 5):
$dF_t=\sigma F_t\,dW^T_t$. A European call with expiry $T$ is then worth
\[
C=P(0,T)\bigl(F\Phi(d_1)-K\Phi(d_2)\bigr),\qquad d_{1,2}=\frac{\ln(F/K)}{\sigma\sqrt T}\pm\tfrac12\sigma\sqrt T .
\]
\end{definition}

The \omterm{thm:dv:black-scholes-three-ways:formula}{Black--Scholes formula} is the special case $F=Se^{(r-q)T}$,
$P(0,T)=e^{-rT}$. Black published the forward version in 1976 for commodity
contracts; the futures options of One Quant Book 3, chapter 12, the caplets of
One Quant Book 2, chapter 13, and the FX options of chapter 20 are all priced
with it, and implied volatility is defined through it (One Quant Book 1,
chapter 25). When the rate is stochastic, discounting with $P(0,T)$ and taking
the expectation under the forward measure keeps the formula exact, provided the
forward itself is lognormal.

\begin{center}
\footnotesize
\begin{tabular}{@{}p{2.4cm}p{4.9cm}p{5.4cm}@{}}
\toprule
Route & What it assumes & What it teaches \\
\midrule
Replication & Itô's formula, continuous hedging & why $\mu$ disappears; the hedge ratio; the equation generalises to any payoff and to numerical grids (chapter 22) \\
Martingale & a measure, a numeraire & which expectation is taken; the choice of numeraire simplifies payoffs; generalises to Monte Carlo (chapter 23) and to other dynamics (chapters 9--13) \\
Binomial limit & two-state steps & that the result is arbitrage, not calculus; American exercise (chapter 6) \\
\bottomrule
\end{tabular}
\end{center}

\begin{omfigure}
\begin{tikzpicture}[font=\small,
  bx/.style={draw=omDef,thick,rounded corners=2pt,fill=omDef!6,align=center,text width=3.3cm,inner sep=4pt},
  res/.style={draw=omThm,very thick,rounded corners=2pt,fill=omThm!8,align=center,text width=3.6cm,inner sep=5pt}]
\node[bx] (a) at (0,2.2) {hold $\partial_SV$ shares:\\Itô removes $dW$};
\node[bx] (b) at (0,0) {Girsanov: $\mathbb Q$, then\\the share numeraire};
\node[bx] (c) at (0,-2.2) {$n$-step tree,\\central limit theorem};
\node[bx] (a2) at (4.6,2.2) {pricing equation\\+ Feynman--Kac};
\node[bx] (b2) at (4.6,0) {$\mathbb Q(S_T>K)=\Phi(d_2)$\\$\mathbb Q^S(S_T>K)=\Phi(d_1)$};
\node[bx] (c2) at (4.6,-2.2) {binomial sums\\$\to$ normal integrals};
\node[res] (f) at (9.3,0) {$Se^{-qT}\Phi(d_1)$\\$-\,Ke^{-rT}\Phi(d_2)$};
\draw[-{Stealth},thick] (a) -- (a2);
\draw[-{Stealth},thick] (b) -- (b2);
\draw[-{Stealth},thick] (c) -- (c2);
\draw[-{Stealth},thick] (a2.east) -- (f.north west);
\draw[-{Stealth},thick] (b2.east) -- (f.west);
\draw[-{Stealth},thick] (c2.east) -- (f.south west);
\end{tikzpicture}

\omcaption{Three routes to one formula: replication (top), the martingale
expectation (middle), the limit of the binomial tree (bottom). Each needs a
different piece of mathematics and each generalises to a different numerical
method.}\label{fig:dv:black-scholes-three-ways:routes}
\end{omfigure}

\begin{method}[Implied volatility, robustly]\label{met:dv:black-scholes-three-ways:iv}
To invert Black's formula at a price $c$:
\begin{enumerate}
\item Check $P(0,T)(F-K)^+\le c<P(0,T)F$ (call); outside these bounds there is no
volatility, and the quote or the forward is wrong.
\item Bracket: the price is increasing in $\sigma$; double an upper bound until it
prices above $c$.
\item Start from a closed-form guess on the time value and run Newton's method
with vega as derivative, falling back to bisection whenever a step leaves the
bracket. Stop on relative price error ($10^{-12}$).
\end{enumerate}
Newton alone fails far from the money, where vega is tiny and a step overshoots;
production libraries use rational approximations that converge in two
iterations to machine precision (Jäckel).
\end{method}

\section{Tutorial: one price, three ways}\label{tut:dv:black-scholes-three-ways:tut}

\begin{tutorial}
\textbf{Goal.} Implement Black's formula and a robust implied-volatility solver,
then price the reference call by the formula, a tree and Monte Carlo, and replicate
it along a path. \textbf{End state:}
\cref{fig:dv:black-scholes-three-ways:replication,fig:dv:black-scholes-three-ways:mc}
and the numbers of the weekend problem.
\begin{tutsteps}
\item \textbf{The formula}, on the forward, with a normal cdf accurate in the tails:
\omcode{code/firm/bs/firm_bs.py}{23}{36}{Black's formula on a forward, and Black--Scholes on a spot.}\label{lst:dv:black-scholes-three-ways:black}
\item \textbf{The solver} of \cref{met:dv:black-scholes-three-ways:iv}:
\omcode{code/firm/bs/firm_bs.py}{61}{94}{Implied volatility: bounds, bracket, guess, safeguarded Newton.}\label{lst:dv:black-scholes-three-ways:iv}
\item \textbf{Monte Carlo}: draw $S_T$ exactly from its lognormal law, average the
discounted payoff, and report the standard error.
\omcode{code/derivatives/03-black-scholes-three-ways/python/dv_blackscholes.py}{24}{34}{Monte Carlo under the risk-neutral measure, with its standard error.}\label{lst:dv:black-scholes-three-ways:mc}
\item \textbf{Run} \texttt{dv\_blackscholes.problem()} and
\texttt{fig\_blackscholes.py}; the C++20 and Rust twins in
\texttt{code/firm/bs/} pass the same round-trip tests.
\end{tutsteps}
\textbf{What to change next.} Turn on antithetic variates and compare the standard
error at equal work; rebalance the \omterm{def:dv:no-arbitrage-and-the-fundamental-theorems:claim}{replicating portfolio} weekly and hourly and
compare the gap at expiry with the daily one.
\end{tutorial}

\begin{omfigure}
\begin{tikzpicture}
\begin{semilogxaxis}[width=10.5cm,height=5cm,
  xlabel={number of paths $M$ (log scale)},ylabel={call price},
  tick label style={font=\small},label style={font=\small},
  xmin=80,xmax=1.3e6,ymin=7.5,ymax=15.5,grid=major,grid style={gray!25},
  legend columns=3,legend style={font=\footnotesize,at={(0.5,-0.3)},anchor=north,draw=none,fill=none,/tikz/every even column/.append style={column sep=8pt}},
  table/col sep=comma]
\addplot[name path=hi,omDef!40,thin] table[x=m,y=hi]{figdata/derivatives/03-black-scholes-three-ways/mc_convergence.csv};
\addplot[name path=lo,omDef!40,thin] table[x=m,y=lo]{figdata/derivatives/03-black-scholes-three-ways/mc_convergence.csv};
\addplot[omDef!20] fill between[of=hi and lo];
\addplot[omDef,very thick,mark=*,mark size=1.5pt] table[x=m,y=estimate]{figdata/derivatives/03-black-scholes-three-ways/mc_convergence.csv};
\addplot[omThm,thick,dashed] table[x=m,y=formula]{figdata/derivatives/03-black-scholes-three-ways/mc_convergence.csv};
\legend{,,two standard errors,Monte Carlo estimate,formula 10.4506}
\end{semilogxaxis}
\end{tikzpicture}

\omcaption{Monte Carlo estimates of the reference call on the first $M$ of one
million draws. The band of two standard errors shrinks like $1/\sqrt M$: a
hundred times more paths buy one more significant digit. Data: the
tutorial.}\label{fig:dv:black-scholes-three-ways:mc}
\end{omfigure}

\section{Build: the Black--Scholes kernels}\label{bld:dv:black-scholes-three-ways:bs}

\begin{build}
\bfield{Purpose} The inner loop of the miniature firm: every surface fit, every
Greek report, the risk engine of One Quant Book 6 and the pricing library of
chapter 28 call these functions millions of times a day, so they exist in
Python, C++20 and Rust with the same tests.
\bfield{Interface} \texttt{black(fwd, strike, t, df, vol, right)};
\texttt{bs(spot, strike, t, r, q, vol, right)}; \texttt{greeks(\ldots)} returning
delta, gamma, vega, theta, rho, vanna and volga;
\texttt{implied\_vol(price, fwd, strike, t, df, right)}. C++:
\texttt{firm::bs} in \texttt{cpp/firm\_bs.hpp}; Rust: crate \texttt{firm\_bs}.
\bfield{Rules} Price on the forward; $\Phi$ through the complementary error
function so that deep out-of-the-money prices keep their relative accuracy; a
price outside the arbitrage bounds raises, never returns a volatility.
\bfield{Acceptance tests} \texttt{code/firm/bs/tests/} and the C++ and Rust
tests: the textbook value 10.450583572185579; parity with a dividend yield; each
Greek against a central difference; implied-volatility round trips to a relative
$10^{-11}$ over five decades of moneyness and maturities from a day to 30 years.
\bfield{Stretch} Replace the solver by a rational-approximation method and
measure iterations and time per inversion in the three languages.
\end{build}

\begin{omsources}
\item F.~Black and M.~Scholes, ``The pricing of options and corporate
liabilities'', \emph{Journal of Political Economy} 81 (1973) 637--654.
\item R.~C.~Merton, ``Theory of rational option pricing'', \emph{Bell Journal of
Economics and Management Science} 4 (1973) 141--183.
\item F.~Black, ``The pricing of commodity contracts'', \emph{Journal of
Financial Economics} 3 (1976) 167--179.
\item P.~Jäckel, ``Let's be rational'', \emph{Wilmott} (2015) 40--53.
\item H.~A.~Baker, ``Cboe at 50'', \emph{Financial History} (Spring 2023).
\end{omsources}

\section{Exercises}

\begin{exercise}[$\star$]\label{exo:dv:black-scholes-three-ways:1}
Price the six-month call and put struck at 100 on a share at 100 with $r=3\%$,
no dividend and $\sigma=25\%$.
\end{exercise}

\begin{exercise}[$\star$]\label{exo:dv:black-scholes-three-ways:2}
Check that $V=S$ (holding the share, with $q=0$) and $V=e^{rt}$ (the bank account)
solve the \omterm{def:dv:black-scholes-three-ways:pde}{Black--Scholes equation}. What does $V=Se^{-q(T-t)}$ represent when
$q>0$?
\end{exercise}

\begin{exercise}[$\star$]\label{exo:dv:black-scholes-three-ways:3}
A futures contract trades at 80. Price a three-month call struck at 85 with a
discount factor of 0.99 and $\sigma=35\%$.
\end{exercise}

\begin{exercise}[$\star\star$]\label{exo:dv:black-scholes-three-ways:4}
A trader expects the share of the reference call to return 15\% a year and
concludes that the call is underpriced at 10.4506. What would the trader have to
believe for the call to be mispriced, and how would the trader profit?
\end{exercise}

\begin{exercise}[$\star\star$]\label{exo:dv:black-scholes-three-ways:5}
For the reference call, give the risk-neutral probability of exercise and the
real-world probability with $\mu=12\%$. Which one does the price use, and why is
neither the hedge ratio?
\end{exercise}

\begin{exercise}[$\star\star$]\label{exo:dv:black-scholes-three-ways:6}
Differentiate the call with respect to the strike to find the price of a digital
call that pays 1 if $S_T>K$, and evaluate it for the reference parameters.
\end{exercise}

\begin{exercise}[$\star\star\star$]\label{exo:dv:black-scholes-three-ways:7}
\emph{Coding.} With \texttt{implied\_vol}, find the volatility at which the
reference call (spot 100, strike 100, one year, $r=5\%$) is worth exactly 10.00.
\end{exercise}

\begin{exercise}[$\star\star\star$]\label{exo:dv:black-scholes-three-ways:8}
\emph{Find the flaw.} ``Our Monte Carlo with 10\,000 paths gives 10.64 for the
reference call, 0.19 above the formula. The formula ignores fat tails: we should
quote 10.64.''
\end{exercise}

\section{Problem: One Price, Three Ways}

\begin{problem}[{Weekend problem --- deriving the desk's reference price three times}]\label{pb:dv:black-scholes-three-ways:1}
The reference call of \cref{ex:dv:black-scholes-three-ways:ref} ($S=K=100$,
$T=1$, $r=5\%$, $q=0$, $\sigma=20\%$) is to be priced by every route of the
chapter, and the answers compared.

\medskip\noindent\textbf{Part I --- Replication.}
\begin{enumerate}
\item Give the delta, gamma and theta (per year) of the call.
\item Check numerically that the call satisfies the \omterm{def:dv:black-scholes-three-ways:pde}{Black--Scholes equation}.
\item Along the chapter's simulated path, what is the gap between the \omterm{def:dv:no-arbitrage-and-the-fundamental-theorems:claim}{replicating
portfolio} and the option at expiry?
\item Why is the gap not zero, and why does its sign not matter for the price?
\item Where did the path's drift of 12\% go?
\end{enumerate}

\medskip\noindent\textbf{Part II --- The martingale route.}
\begin{enumerate}[resume]
\item Give $d_1$, $d_2$, $\Phi(d_1)$ and $\Phi(d_2)$.
\item Give the call and the put by the formula, and check parity.
\item Give the risk-neutral and the real-world ($\mu=12\%$) probabilities of exercise.
\item Give the digital call's price.
\item Compare the call with the at-the-money approximation $0.4S\sigma\sqrt T$; why
does it understate here?
\end{enumerate}

\medskip\noindent\textbf{Part III --- Numerical routes.}
\begin{enumerate}[resume]
\item Give the 1\,000-step Cox--Ross--Rubinstein price and its error.
\item Give the Monte Carlo price with $10^6$ paths, its standard error and its error.
\item How many paths would reach the tree's accuracy?
\item Give the Monte Carlo price with $10^4$ paths and its standard error.
\item Which method would you use for 10\,000 European calls, and which for one
American put?
\end{enumerate}

\medskip\noindent\textbf{Part IV --- Judgement.}
\begin{enumerate}[resume]
\item Which assumption of the model fails most visibly in the market?
\item Why do desks still quote in Black--Scholes volatility?
\item What would change in the formula if the share paid a 2\% dividend yield?
\item State the \emph{named result}: the reference price by the three routes, with the
error of each.
\item In one sentence: why does the drift not appear?
\end{enumerate}
\end{problem}

\section{Interview questions}

\begin{interviewq}[$\star$ \iqroles{researcher, trader}]\label{iq:dv:black-scholes-three-ways:1}
Derive the \omterm{def:dv:black-scholes-three-ways:pde}{Black--Scholes equation}.
\end{interviewq}

\begin{interviewq}[$\star$ \iqroles{trader}]\label{iq:dv:black-scholes-three-ways:2}
A one-year at-the-money call on a share at 200 with 30\% volatility: roughly
what is it worth? (Ignore rates.)
\end{interviewq}

\begin{interviewq}[$\star\star$ \iqroles{researcher, bank}]\label{iq:dv:black-scholes-three-ways:3}
What are $\Phi(d_1)$ and $\Phi(d_2)$?
\end{interviewq}

\begin{interviewq}[$\star\star$ \iqroles{researcher}]\label{iq:dv:black-scholes-three-ways:4}
Why does the expected return of the share not appear in the formula? Would
it appear if you could not trade the share?
\end{interviewq}

\begin{interviewq}[$\star\star$ \iqroles{trader, risk}]\label{iq:dv:black-scholes-three-ways:5}
List the assumptions of the model. Which failure costs a volatility desk most?
\end{interviewq}

\begin{interviewq}[$\star\star\star$ \iqroles{developer}]\label{iq:dv:black-scholes-three-ways:6}
You must compute implied volatilities for 50 million option quotes a day,
including deep out-of-the-money and nearly expired ones. How do you make the
solver robust and fast?
\end{interviewq}
